\section{Pruebas y Evidencias}

\subsection{Pruebas Unitarias y de Integración}
Siguiendo la metodología guiada por pruebas (TDD), se implementaron pruebas unitarias sobre los casos de uso principales utilizando dobles de prueba (\textit{fakes}) para los puertos del hipervisor y del DNS. Esto garantiza que la lógica de orquestación y las rutinas de compensación ante fallos se validen de forma continua e independiente de la infraestructura física.

\subsection{Evidencias de Despliegue en Entorno Real}
A continuación se relacionan las evidencias del ciclo de vida de aprovisionamiento de las instancias solicitadas en la rúbrica del proyecto (\texttt{paisajes} y \texttt{negocios}):

\begin{figure}[htbp]
\centering
\includegraphics[width=0.9\linewidth]{captura-dashboard-vacio.png}
\caption{Panel de gestión inicial en estado vacío.}
\label{fig:dashboard_vacio}
\end{figure}

\begin{figure}[htbp]
\centering
\includegraphics[width=0.9\linewidth]{captura-formulario-subida.png}
\caption{Carga del archivo comprimido para la instancia \texttt{paisajes}.}
\label{fig:form_subida}
\end{figure}

\begin{figure}[htbp]
\centering
\includegraphics[width=0.9\linewidth]{captura-instancia-paisajes.png}
\caption{Instancia \texttt{paisajes} desplegada y visible en el panel.}
\label{fig:instancia_paisajes}
\end{figure}

\begin{figure}[htbp]
\centering
\includegraphics[width=0.9\linewidth]{captura-instancia-negocios.png}
\caption{Instancia \texttt{negocios} desplegada concurrentemente.}
\label{fig:instancia_negocios}
\end{figure}

